\section{The \texttt{Source\_pt} McXtrace Component}
Release: McXtrace 0.1

An x-ray point source

\subsection*{Identification}
\begin{itemize}
  \item \textbf{Author:} Erik Knudsen
  \item \textbf{Origin:} Risoe
  \item \textbf{Date:} June 29th, 2009
\end{itemize}

\subsection*{Description}
A simple source model emitting photons from a point source uniformly into 4pi. A square target centered on the Z-axis restricts the beam to that aperture. If an input spectrum datafile (spectrum\_file) is not specified, the beam is restricted to emit photons between E0+-dE keV, or lambda0+-dlambda \AA{}, whichever is given. The input spectrum file should be formatted such that x-ray energy/wavelength is in the first column and the intensity in the second. Any preceding lines starting with \# are considered part of the file header. If a datafile is given, a nonzero E¤0 value indicates that is is parametrized by energy ( in keV) as opposed to wavelength (in \AA{}). Wavelength is the default. Flux is given in the unit photons/s

Example: Source\_pt(dist=1,focus\_xw=0.1,focus\_yh=0.1, lamda=0.231, dlambda=0.002)

\subsection*{Input parameters}
Parameters in \textbf{boldface} are required; the others are optional.

\begin{longtable}{p{0.22\textwidth}p{0.12\textwidth}p{0.46\textwidth}p{0.14\textwidth}}
\toprule
\textbf{Name} & \textbf{Unit} & \textbf{Description} & \textbf{Default} \\
\midrule
\endhead
focus\_xw & m & Width of target & 0 \\
focus\_yh & m & Height of target & 0 \\
focus\_x0 & m & x-cocordinate of target centre. & 0 \\
focus\_y0 & m & y-coordinate of target centre. & 0 \\
flux & ph/s & Total flux radiated from the source. & 0 \\
dist & m & Distance from source plane to sampling window. & 1 \\
E0 & keV & Mean energy of xrays. & 0 \\
dE & keV & Energy half spread of x-rays. & 0 \\
lambda0 & \AA{} & Mean wavelength of x-rays. & 0 \\
dlambda & \AA{} & Wavelength half spread of x-rays (flat or gaussian sigma). & 0 \\
phase & rad & Set phase to something given. & 0 \\
randomphase & 0/1 & If nonzero, the phase of the emotted photon is random, i.e. source is fully incoherent. otherwise the value of phase is used. & 1 \\
gauss & 1 & Gaussian (1) or Flat (0) energy/wavelength distribution & 0 \\
spectrum\_file & string & File from which to read an input spectrum. & "" \\
verbose & 1 & Output more information runtime. & 0 \\
\bottomrule
\end{longtable}

\subsection*{Links}
\begin{itemize}
  \item Component source code found in file \texttt{Source\_pt.comp}.
\end{itemize}
\IfFileExists{sources/Source_pt_static.tex}{\input{sources/Source_pt_static.tex}}{}